%! Representing a Categorical Variable — Unit 1, Lesson 1.2 — Guided Notes & Practice
% THE in-class document, in the MAIN IDEAS / NOTES shape (see LESSON_SHAPE.md §1): the
% vocabulary box, then ONE two-column guidednotes table. Each row is one idea — a label on the
% left; on the right a terse definition the student completes, then a grid of short numbered
% problems with reserved work space. Problems are numbered 1..N through the whole component.
%   I DO   : the teacher fills the definition lines and works the FIRST problem of each grid.
%   WE DO  : the class works the rest of each grid, students holding the pen; the last row,
%            Guided Practice, is worked entirely together in a fresh context.
%   YOU DO : nothing on this page — ../homework/ IS the individual practice, started in the
%            last minutes of the period. No independent set, no reflection box, no objectives
%            box (the targets are on the cover).
% Vocabulary is GIVEN here, not discovered. The crux is the WINS FIRST row — South has more
% walkers AND North has the higher walking rate, both true, and the bigger group wins on every
% raw count before anyone counts anything (EK UNC-1.E.1). The plurality trap (problem 6) and
% the wrong-total trap (problem 19) are problems in a grid, not callouts.
%
% This lesson absorbs TWO CED topics (1.3 tables and 1.4 graphs), so the Bar Graphs row keeps
% the mechanics to one pre-drawn chart and rides the axis idea into the crux row.
%
% THE DATA: North High, all 400 students — Bus 180, Car 120, Walk 80, Bike 20.
%           South High, all 1,000 students — Walk 150, Bike 20 (the district's side-by-side chart).
%           Two Shifts (guided practice) — Morning 48/108/84 of 240; Evening 18/30/32 of 80.
%
% PAGE LOCKSTEP with notes_key/: every \blank{} here becomes a short \ans{} there, every
% \termblank becomes a \termans, and every \pcell{n}{...}{H}{} gets its answer in the fourth
% argument. Heights and blank widths are sized to the answers so the two files set identically.
% A guidednotes row cannot break across pages — a long idea is split into sub-rows (a `\\` with
% no \hline, then `& ...`) so the table can break between them.
% NO SKETCH QUESTIONS: both displays are pre-drawn in TikZ and read, never constructed.
\documentclass[12pt]{article}
\usepackage{apstats-article}
\usepackage{apstats-boxes}

\begin{document}

\pageheader{Unit 1, Lesson 1.2}{Guided Notes \& Practice}
% NAMESTRIP: no \namedateperiod here — the name row lives on the cover only.

\vspace{2pt}

\boxguard[16]
\begin{vocabbox}
\small
Fill in each term \textbf{as we name it} in the notes below.
\par
\termblank{Frequency table}
\par
\termblank{Relative frequency table}
\par
\termblank{Bar graph}
\par
\termblank{The comparison rule}
\par
\end{vocabbox}

\small
\begin{guidednotes}
% ── ROW 1: frequency and relative frequency tables ───────────────────────────
\mainidea[Counts \&]{Shares} &
A categorical column cannot be averaged, but it can be \textbf{counted}. A \textbf{frequency
table} lists how many cases fall in each category; the counts add to the \blank{2.2cm}.
A \textbf{relative frequency table} lists the share of the group in each category; the shares
always add to \blank{2.3cm}, because every case is counted in exactly \blank{1.1cm} category.

\vspace{2pt}
\textbf{1.} North High asked all \textbf{400} of its students how they get to school. Fill the
shares:

\vspace{1pt}
\renewcommand{\arraystretch}{1.1}
\begin{tabularx}{\linewidth}{@{} l Y Y @{}}
\rowcolor{sky}
\textbf{How they get to school} & \textbf{Frequency} (count) & \textbf{Relative frequency} \\
\hline
Bus  & 180 & \blank{2.4cm} \\
Car  & 120 & \blank{2.4cm} \\
Walk &  80 & \blank{2.4cm} \\
Bike &  20 & \blank{2.4cm} \\
\hline
\textbf{Total} & \textbf{400} & \blank{2.4cm} \\
\end{tabularx}
\notesprompt{Read your table.}
\begin{probgrid}{|Y|Y|Y|}
\pcell{2}{What do the four counts add to, and why must they?}{1.2cm}{} &
\pcell{3}{What do the four shares add to?}{1.2cm}{} &
\pcell{4}{Eighty of 400 walk. Write that share as a fraction, a decimal, a percent, and a
rate.}{1.2cm}{} \\ \hline
\end{probgrid}
\\ \hline
% ── ROW 2: plurality vs majority, and the context habit ──────────────────────
\mainidea[``Most'' is a promise]{About half} &
The largest single category is a \blank{1.9cm}. A \textbf{majority} needs more than
\blank{1.2cm} of the group --- not the same thing. On the AP exam a share always comes with
its \blank{1.6cm}: not ``20\%,'' but \emph{20\% \textbf{of North students} walk to school}.
\notesprompt{Check the claim.}
\begin{probgrid}{|Y|Y|Y|}
\pcell{5}{A flyer says ``Most North students take the bus.'' What share takes the bus?}{1.2cm}{} &
\pcell{6}{Is that a majority? Say why or why not.}{1.2cm}{} &
\pcell{7}{Rewrite ``20\%'' as a sentence an AP reader would accept.}{1.2cm}{} \\ \hline
\end{probgrid}
\\ \hline
% ── ROW 3: bar graphs ────────────────────────────────────────────────────────
\mainidea{Bar graphs} &
The same North High column, drawn. Each bar stands for one \blank{1.9cm}, and its height is
the \blank{1.3cm} in that category --- or the \blank{2.9cm}, if the vertical axis is relabelled.
The bars have \blank{1.2cm} between them: categories are separate groups, not points on a
number line, so the order of the bars carries no meaning.

\begin{center}
\begin{tikzpicture}[scale=0.5]
  \draw[linegray, very thin] (0,1) -- (5.4,1);
  \draw[linegray, very thin] (0,2) -- (5.4,2);
  \draw[linegray, very thin] (0,3) -- (5.4,3);
  \draw[->, thick] (0,0) -- (0,4.1);
  \draw[->, thick] (0,0) -- (5.6,0);
  \foreach \y/\lab in {1/50, 2/100, 3/150} {
    \draw (-0.08,\y) -- (0.08,\y) node[left, font=\tiny, xshift=-2pt] {\lab};}
  \foreach \x/\h/\lab in {0.5/3.6/Bus, 1.8/2.4/Car, 3.1/1.6/Walk, 4.4/0.4/Bike} {
    \fill[navy] (\x,0) rectangle ++(0.8,\h);
    \node[below, font=\tiny] at (\x+0.4,-0.05) {\lab};}
  \node[rotate=90, font=\tiny] at (-0.75,2.0) {Students};
\end{tikzpicture}
\end{center}
\\
& \notesprompt{Read the graph.}
\begin{probgrid}{|Y|Y|Y|}
\pcell{8}{How tall is the Walk bar, and what is it counting?}{1.0cm}{} &
\pcell{9}{What share of North is that?}{1.0cm}{} &
\pcell{10}{Add the four bar heights. What do you get?}{1.0cm}{} \\ \hline
\pcell{11}{Why are there gaps between the bars?}{1.3cm}{} &
\pcell{12}{Relabel the axis 0.45, 0.30, 0.20, 0.05. Which bars change height?}{1.3cm}{} &
\pcell{13}{Before reading any bar graph, what is the first question to ask?}{1.3cm}{} \\ \hline
\end{probgrid}
\\ \hline
% ── ROW 4: the comparison rule — THE CRUX ────────────────────────────────────
\mainidea[The bigger group]{Wins first} &
South High is in the same district and surveyed all \textbf{1{,}000} of its students. The
district posted the chart on the right to compare the two schools.

\vspace{2pt}
\begin{minipage}[t]{0.46\linewidth}
\renewcommand{\arraystretch}{1.15}
\begin{tabularx}{\linewidth}{@{} l Y Y Y @{}}
\rowcolor{sky}
\textbf{School} & \textbf{Walk} & \textbf{Bike} & \textbf{Total} \\
\hline
North &  80 & 20 &   400 \\
South & 150 & 20 & 1{,}000 \\
\hline
\end{tabularx}
\end{minipage}
\hfill
\begin{minipage}[t]{0.50\linewidth}
\vspace{-6pt}
\begin{center}
\begin{tikzpicture}[scale=0.5]
  \draw[linegray, very thin] (0,1) -- (6.6,1);
  \draw[linegray, very thin] (0,2) -- (6.6,2);
  \draw[linegray, very thin] (0,3) -- (6.6,3);
  \draw[->, thick] (0,0) -- (0,3.8);
  \draw[->, thick] (0,0) -- (6.8,0);
  \foreach \y/\lab in {1/200, 2/400, 3/600} {
    \draw (-0.08,\y) -- (0.08,\y) node[left, font=\tiny, xshift=-2pt] {\lab};}
  \foreach \x/\hn/\hs/\lab in {0.4/0.9/2.5/Bus, 2.0/0.6/1.65/Car, 3.6/0.4/0.75/Walk,
                               5.2/0.1/0.1/Bike} {
    \fill[navy]   (\x,0)     rectangle ++(0.55,\hn);
    \fill[goldacc](\x+0.6,0) rectangle ++(0.55,\hs);
    \node[below, font=\tiny] at (\x+0.57,-0.05) {\lab};}
  \fill[navy] (5.9,3.3) rectangle ++(0.28,0.28);
  \node[right, font=\tiny] at (6.2,3.44) {North};
  \fill[goldacc] (5.9,2.9) rectangle ++(0.28,0.28);
  \node[right, font=\tiny] at (6.2,3.04) {South};
\end{tikzpicture}
\end{center}
\end{minipage}
\\
& South's bar is taller in almost every category --- and that was guaranteed before anyone
counted anything, because South surveyed \blank{1.9cm} as many students.
\textbf{The comparison rule:} to compare groups of different sizes, use \blank{3.6cm},
never raw counts. But ``how many students in this district walk?'' asks about \blank{1.9cm},
and there a count is exactly the right number --- \emph{say which number you used, and why}.
\notesprompt{Compare the schools.}
\begin{probgrid}{|Y|Y|}
\pcell{14}{Work both walking rates.}{1.3cm}{} &
\pcell{15}{Which school has more walkers? Which has the higher walking rate?}{1.3cm}{} \\ \hline
\pcell{16}{So which school ``walks more''? Answer honestly.}{1.2cm}{} &
\pcell{17}{What one change would make the district's chart a fair comparison?}{1.2cm}{} \\ \hline
\pcell{18}{``How many students in this district walk to school?'' Which number, and why?}{1.3cm}{} &
\pcell{19}{A student divides South's 150 walkers by 1{,}400. What did she compute?}{1.3cm}{} \\ \hline
\end{probgrid}
\\ \hline
% ── ROW 5: guided practice — WE DO, together, fresh context ──────────────────
\mainidea[Guided practice]{Two Shifts} &
A coffee shop records the size of every drink it sells. Morning shift: \textbf{240} orders.
Evening shift: \textbf{80}.

\vspace{1pt}
\renewcommand{\arraystretch}{1.1}
\begin{tabularx}{\linewidth}{@{} l Y Y Y Y @{}}
\rowcolor{sky}
\textbf{Shift} & \textbf{Small} & \textbf{Medium} & \textbf{Large} & \textbf{Total} \\
\hline
Morning &  48 & 108 & 84 & 240 \\
Evening &  18 &  30 & 32 &  80 \\
\hline
\end{tabularx}
\notesprompt{We work these together.}
\begin{probgrid}{|Y|Y|Y|}
\pcell{20}{Which shift sold more large drinks?}{1.0cm}{} &
\pcell{21}{Morning's large-drink rate.}{1.0cm}{} &
\pcell{22}{Evening's large-drink rate.}{1.0cm}{} \\ \hline
\pcell{23}{Which shift sells large drinks at the higher rate? Is that the same question as
problem~20?}{1.4cm}{} &
\pcell{24}{One sentence for the staff newsletter comparing what the shifts sell. Which kind of
number did you use, and why?}{1.4cm}{} &
\pcell{25}{The owner asks, ``Which shift is \emph{busier}?'' Which column answers that --- and
is the answer a rate or a count?}{1.4cm}{} \\ \hline
\end{probgrid}
\\ \hline
\end{guidednotes}

% ── THE NOTES END AT GUIDED PRACTICE ─────────────────────────────────────────
% There is deliberately no "you do" here: ../homework/ is the individual practice,
% started in class in the last 8 minutes while the teacher circulates.

\end{document}
